% !TeX root = ../MQF_Manuale_Master.tex
% Capitoli/MQF_Cap_14_Programmazione_Stocastica_I.tex

\chapter{Programmazione stocastica a due stadi}

\label{cap:programmazione-stocastica-due-stadi}

\medskip

\noindent
\textbf{Obiettivi della lezione}

Il Capitolo 12 ha formulato un problema di Asset--Liability Management
deterministico nel quale la composizione iniziale dell'attivo, le giacenze
di liquidit\`{a} e il funding vengono coordinati lungo pi\`{u} date. La
dimensione temporale \`{e} esplicita, ma l'intera traiettoria dei coefficienti
futuri \`{e} considerata nota al tempo $0$. In queste condizioni il modello
pu\`{o} determinare fin dall'inizio un unico piano completo.

Il presente capitolo abbandona questa ipotesi. Il problema non consiste
semplicemente nel sostituire alcuni coefficienti deterministici con
quantit\`{a} aleatorie. Quando il futuro non \`{e} noto, occorre specificare
anche \emph{quando} le decisioni vengono assunte, \emph{quale informazione}
sia disponibile in quel momento e \emph{quali decisioni} possano essere
modificate dopo avere osservato l'esito dell'incertezza.

Questa struttura conduce alla programmazione stocastica a due stadi. Una
prima decisione viene assunta prima della rivelazione dell'incertezza; una
seconda decisione pu\`{o} invece reagire all'informazione successivamente
osservata. La soluzione non \`{e} quindi costituita soltanto da una scelta
iniziale, ma comprende anche una politica di risposta contingente ai
possibili esiti futuri.

Il capitolo sviluppa questa metodologia mantenendo in parallelo due casi
guida, scelti perch\'e permettono di osservare la stessa struttura matematica
in contesti finanziari ed economici differenti.

Il primo caso prosegue il problema SVB sviluppato nei Capitoli 11--12. Nel
modello deterministico la banca costruisce il proprio piano conoscendo
l'intera traiettoria dei dati futuri. La formulazione stocastica elimina
questa ipotesi: alcune decisioni devono essere assunte prima di conoscere
l'evoluzione delle condizioni economico-finanziarie, mentre altre decisioni
di liquidit\`a e funding possono essere adattate all'informazione
successivamente osservata.

Il secondo caso \`e ispirato al complesso industriale Alcoa di San
Cipri\'an, situato in Galizia, nella provincia di Lugo, tra i comuni di
Cervo e Xove. Il sito, operativo dal 1980, comprende una raffineria di
allumina e un impianto per la produzione di alluminio primario. La produzione
di alluminio mediante elettrolisi rende il costo e la disponibilit\`a
dell'energia elementi economicamente rilevanti per l'attivit\`a produttiva.
Negli anni recenti proprio le condizioni di approvvigionamento energetico
hanno avuto un ruolo centrale nel dibattito sulla continuit\`a e sulla
competitivit\`a dell'impianto.

Il caso utilizzato nel capitolo \`e una rappresentazione stilizzata e non
una ricostruzione dei contratti effettivamente sottoscritti da Alcoa. Il suo
interesse risiede nella struttura della decisione. Una parte
dell'approvvigionamento pu\`o essere definita anticipatamente mediante
decisioni contrattuali, mentre quantit\`a e costi correttivi possono essere
determinati dopo l'osservazione delle condizioni future. Il problema combina
quindi rischio di prezzo, decisioni intertemporali e flessibilit\`a
contrattuale.

Dal punto di vista matematico, il caso \`e particolarmente utile perch\'e
l'incertezza riguarda grandezze che evolvono nel tempo e che possono essere
dipendenti tra loro. Esso consente quindi di collegare i processi stocastici
e la simulazione di traiettorie studiati nei capitoli precedenti alla
costruzione di scenari per un problema di ottimizzazione. Il confronto con
il caso SVB permette inoltre di mostrare due modi differenti di ottenere una
rappresentazione finita dell'incertezza: a partire da un modello
probabilistico discreto nel primo caso e mediante simulazione e successiva
riduzione degli scenari nel secondo.

Al termine del capitolo lo studente dovr\`a essere in grado di:

\begin{itemize}

\item riconoscere la struttura temporale e informativa di un problema
decisionale sotto incertezza, distinguendo decisioni preventive e decisioni
di recourse;

\item rappresentare i dati aleatori mediante processi stocastici e scenari
finiti, comprendendo la relazione tra modello probabilistico e
rappresentazione utilizzata nell'ottimizzazione;

\item formulare e interpretare un problema di programmazione stocastica a
due stadi, dalla funzione di recourse alla forma estesa per scenari;

\item valutare il ruolo economico della flessibilit\`a decisionale e
dell'informazione attraverso i problemi stochastic programming,
wait-and-see ed expected value e gli indicatori EVPI e VSS;

\item applicare il formalismo ai casi SVB--ALM e Alcoa, riconoscendo al
tempo stesso il limite della struttura a due stadi quando l'informazione
viene rivelata progressivamente.

\end{itemize}

\section{Decisioni nel tempo sotto incertezza}

\label{sec:cap14-scegliere-reagire}

Il modello di Asset--Liability Management del Capitolo 12 \`e
multiperiodale: cash flow, fabbisogni di liquidit\`a, giacenze e funding
sono distribuiti lungo un insieme di date

\[
\mathcal T=\{0,1,\ldots,T\}.
\]

La presenza di pi\`u date non rende tuttavia il problema stocastico. Nel
modello deterministico i dati esogeni relativi alle date future sono
considerati noti quando il problema viene risolto al tempo iniziale. Di
conseguenza, l'intero piano pu\`o essere determinato utilizzando fin
dall'inizio informazioni relative a tutto l'orizzonte.

Il passaggio alla programmazione stocastica richiede di modificare questa
ipotesi. Occorre anzitutto rappresentare matematicamente l'evoluzione
temporale dei dati incerti e, successivamente, stabilire quale parte di tale
informazione sia disponibile quando vengono assunte le diverse decisioni.

\subsection{Dalla traiettoria deterministica al processo dei dati incerti}

Supponiamo che alla data $t$ il modello utilizzi $m$ parametri esogeni.
Nel problema deterministico essi possono essere raccolti nel vettore

\[
\bar\xi_t
=
\begin{pmatrix}
\bar\xi_{t1}\\
\vdots\\
\bar\xi_{tm}
\end{pmatrix}
\in\R^m,
\qquad
t\in\mathcal T.
\]

La successione

\[
\{\bar\xi_t\}_{t\in\mathcal T}
\]

descrive quindi una traiettoria deterministica dei dati esogeni lungo
l'orizzonte. La barra sovrapposta indica che tali valori sono trattati come
deterministici e fissati; non indica un valore atteso.

Quando il problema viene risolto al tempo $0$, l'intera traiettoria

\[
\bar\xi_0,\bar\xi_1,\ldots,\bar\xi_T
\]

\`e assunta nota.

Nel caso ALM del Capitolo 12 questa ipotesi permette di determinare
congiuntamente le variabili del piano,

\[
x,
\qquad
g_1,\ldots,g_T,
\qquad
u_1,\ldots,u_{T-1}.
\]

Una variabile come $u_3$ riguarda una decisione di funding riferita alla
data $3$, ma il suo valore pu\`o essere determinato quando il problema viene
risolto al tempo $0$, poich\'e i dati necessari per calcolarlo sono gi\`a
noti. La data alla quale una decisione produce i propri effetti economici
non coincide quindi necessariamente con il momento nel quale l'informazione
utilizzata per determinarla diventa disponibile.

Abbandoniamo ora l'ipotesi che la traiettoria dei dati esogeni sia nota.
Alla data $t$ gli stessi $m$ parametri vengono rappresentati dal vettore
aleatorio

\[
\xi_t
=
\begin{pmatrix}
\xi_{t1}\\
\vdots\\
\xi_{tm}
\end{pmatrix},
\]

dove

\[
\xi_t:(\Omega,\F)\longrightarrow
(\R^m,\B(\R^m)).
\]

La famiglia

\[
\{\xi_t\}_{t\in\mathcal T}
\]

costituisce pertanto un processo stocastico vettoriale. Equivalentemente,
l'intero processo pu\`o essere letto attraverso l'applicazione

\[
(t,\omega)
\longmapsto
\xi_t(\omega),
\qquad
(t,\omega)\in\mathcal T\times\Omega.
\]

Questa rappresentazione riprende direttamente il formalismo dei processi
stocastici introdotto nel Capitolo 5, estendendo lo spazio degli stati dal
caso scalare al caso vettoriale.

Per ogni $t$ fissato,

\[
\omega\longmapsto\xi_t(\omega)
\]

\`e un vettore aleatorio che descrive congiuntamente i parametri incerti
riferiti alla data $t$. Per ogni $\omega$ fissato, invece,

\[
t\longmapsto\xi_t(\omega)
\]

descrive una traiettoria del processo, cio\`e una possibile evoluzione
congiunta dei dati esogeni lungo il tempo.

La natura vettoriale di $\xi_t$ consente di rappresentare simultaneamente
pi\`u fonti di incertezza alla stessa data, mentre la successione temporale

\[
\xi_0,\xi_1,\ldots,\xi_T
\]

consente di rappresentarne l'evoluzione. Nessuna ipotesi di indipendenza
tra le componenti di $\xi_t$, n\'e tra valori assunti a date differenti,
\`e implicata da questa notazione: tali dipendenze fanno parte della legge
congiunta del processo.

\subsection{Dal processo stocastico al vettore dei dati aleatori}

La programmazione stocastica a due stadi utilizza normalmente una notazione
pi\`u compatta. I dati incerti che diventano rilevanti tra il primo e il
secondo stadio vengono raccolti in un unico vettore aleatorio, indicato con

\[
\xi.
\]

Quando il problema utilizza l'intera evoluzione futura del processo e i dati
al tempo $0$ sono gi\`a osservati, tale vettore pu\`o essere costruito
impilando i valori futuri:

\[
\xi
=
\begin{pmatrix}
\xi_1\\
\vdots\\
\xi_T
\end{pmatrix}
=
\begin{pmatrix}
\xi_1' & \cdots & \xi_T'
\end{pmatrix}',
\]

e quindi

\[
\xi:\mathcal T \times\Omega\longrightarrow\Xi\subseteq\R^{m}.
\]

Per ogni $\omega\in\Omega$,

\[
\xi(\omega)
=
\begin{pmatrix}
\xi_1(\omega)\\
\vdots\\
\xi_T(\omega)
\end{pmatrix}
\]

rappresenta una possibile traiettoria congiunta dei dati incerti, espressa
come un unico vettore.

Non tutti i problemi richiedono necessariamente l'intera traiettoria. Se
soltanto alcuni valori del processo sono rilevanti per il secondo stadio,
$\xi$ raccoglie soltanto tali dati. Il simbolo $\xi$ indica quindi il
vettore aleatorio dei dati che entrano nella formulazione stocastica,
mentre $\{\xi_t\}_{t\in\mathcal T}$ descrive il processo temporale dal quale
tali dati hanno origine.

Questa distinzione sar\`a utile nei due casi guida del capitolo. Nel caso
SVB--ALM la distribuzione dei dati del secondo stadio verr\`a costruita a
partire dall'evoluzione dello stato economico-finanziario. Nel caso di
approvvigionamento energetico, invece, una realizzazione di $\xi$ potr\`a
raccogliere una traiettoria di dati simulati lungo pi\`u date.

\subsection{Informazione disponibile e possibilit\`a di reazione}

La specificazione del processo

\[
\{\xi_t\}_{t\in\mathcal T}
\]

descrive come possono evolvere i dati esogeni. Essa non determina per\`o,
da sola, la struttura del problema decisionale. Occorre stabilire quali
informazioni siano disponibili nel momento in cui una decisione viene
assunta.

Consideriamo una decisione iniziale

\[
x
\]

che deve essere scelta prima che siano osservati i dati aleatori raccolti
in $\xi$. Poich\'e la realizzazione futura non \`e ancora nota, $x$ non
pu\`o essere scelta separatamente per ciascuna possibile realizzazione
$\xi(\omega)$.

Supponiamo che, dopo la rivelazione dell'informazione rappresentata da
$\xi$, sia possibile assumere una seconda decisione. Poich\'e questa
decisione viene presa dopo l'osservazione dell'esito, essa pu\`o dipendere
dalla realizzazione osservata e viene indicata con

\[
y(\xi).
\]

La struttura essenziale del problema a due stadi pu\`o quindi essere
rappresentata come

\[
x
\longrightarrow
\xi
\longrightarrow
y(\xi).
\]

\begin{definition}[Decisione di primo stadio e decisione di recourse]
La decisione $x$, che deve essere assunta prima della rivelazione
dell'informazione rappresentata da $\xi$, si dice \emph{decisione di primo
stadio}.

La decisione $y(\xi)$, assunta dopo tale rivelazione e quindi adattabile
alla realizzazione osservata, si dice \emph{decisione di secondo stadio} o
\emph{decisione di recourse}.
\end{definition}

Il termine \emph{recourse} non identifica una specifica operazione
economica. Esprime la possibilit\`a di reagire alle condizioni che si sono
effettivamente manifestate dopo che la decisione iniziale \`e stata fissata.

La legge del processo stocastico e la struttura delle decisioni svolgono
pertanto funzioni differenti. La prima descrive l'incertezza; la seconda
specifica quali decisioni devono essere fissate prima che tale incertezza
sia osservata e quali possano invece utilizzare l'informazione
successivamente disponibile.

\begin{remark}[Date temporali e stadi decisionali]
L'indice $t$ di $\xi_t$ identifica la data alla quale si riferiscono i dati
aleatori. Esso non identifica automaticamente uno stadio decisionale.

Un problema pu\`o contenere dati riferiti a molte date,

\[
\xi_1,\ldots,\xi_T,
\]

e avere tuttavia soltanto due stadi decisionali. Ci\`o accade quando, ai
fini del modello, i dati raccolti nel vettore $\xi$ vengono rivelati nello
stesso passaggio informativo prima delle decisioni di secondo stadio.

La distinzione tra periodo temporale e stadio decisionale dipende quindi
dalla struttura dell'informazione, non dal numero di indici temporali
presenti nel modello.
\end{remark}

Questa osservazione individua il punto nel quale la programmazione stocastica
si collega alla teoria dell'informazione sviluppata nei capitoli precedenti.
Dire che una decisione viene assunta ``prima'' o ``dopo'' l'osservazione di
$\xi$ deve essere tradotto matematicamente specificando rispetto a quale
sigma-algebra essa sia misurabile. Tale formalizzazione verr\`a sviluppata
nella sezione successiva mediante la filtrazione.

\subsection{Caso SVB--ALM}

La struttura astratta

\[
x
\longrightarrow
\xi
\longrightarrow
y(\xi)
\]

pu\`o ora essere letta utilizzando direttamente le variabili dei due casi
guida. In questo modo la distinzione tra decisioni preventive, dati incerti
e decisioni di recourse assume un contenuto economico preciso.

Nel caso SVB la composizione iniziale dell'attivo rimane quella introdotta
nel Capitolo 11:

\[
x=
\begin{pmatrix}
x_1\\
x_2\\
x_3\\
x_4
\end{pmatrix},
\]

dove

\[
\begin{array}{ll}
x_1 & \text{cassa e riserve},\\
x_2 & \text{titoli a breve scadenza},\\
x_3 & \text{titoli a lunga scadenza},\\
x_4 & \text{prestiti e altri impieghi meno liquidi}.
\end{array}
\]

Queste quattro variabili sono scelte al tempo $0$ e soddisfano il vincolo

\[
x_1+x_2+x_3+x_4=A_0.
\]

Il passaggio al modello ALM del Capitolo 12 non modifica la natura di $x$.
Introduce invece quattro date future e la matrice

\[
A^{\mathrm{CF}}
=
(a_{tj}),
\qquad
t=1,\ldots,4,
\qquad
j=1,\ldots,4,
\]

che determina la distribuzione temporale dei cash flow prodotti dalle quattro
classi di attivo.

Alle diverse date la banca deve soddisfare i fabbisogni

\[
d_1,d_2,d_3,d_4.
\]

Le decisioni che consentono di gestire la liquidit\`a nel corso
dell'orizzonte sono

\[
g_1,g_2,g_3,g_4,
\]

dove $g_t$ rappresenta la giacenza di liquidit\`a dopo il soddisfacimento
del fabbisogno alla data $t$, e

\[
u_1,u_2,u_3,
\]

dove $u_t$ rappresenta il funding raccolto alla data $t$ e rimborsato alla
data successiva.

Nel Capitolo 12 l'ambiente economico-finanziario \`e rappresentato dalla
traiettoria deterministica

\[
\bar Z_t
=
(\bar R_t,\bar C_t,\bar F_t),
\]

considerata interamente nota al tempo $0$. Nel passaggio al modello
stocastico questa ipotesi viene modificata esplicitamente: l'ambiente esterno
\`e descritto dal processo

\[
Z_t
=
(Z_t^R,Z_t^C,Z_t^F),
\]

la cui evoluzione futura non \`e nota quando viene scelta la composizione
iniziale $x$.

In particolare, il fabbisogno di liquidit\`a viene ora fatto dipendere dallo
stato esterno:

\[
d_t=d(Z_t).
\]

Si tratta di una modifica rispetto al modello deterministico, nel quale
$d_t$ era un coefficiente noto al tempo $0$. Le altre relazioni del caso
SVB--ALM conservano invece il significato attribuito loro nel Capitolo 12
finch\'e non venga esplicitamente dichiarata una diversa specificazione.

La struttura decisionale diventa quindi

\[
\begin{array}{ccccc}
x
&
\longrightarrow
&
Z_t
&
\longrightarrow
&
g_t,\;u_t,
\end{array}
\]

con

\[
d_t=d(Z_t).
\]

Il vettore $x$ deve essere fissato senza conoscere gli stati futuri. Le
variabili $g_t$ e $u_t$, che nel modello deterministico potevano essere
calcolate tutte al tempo $0$, acquistano nel modello stocastico il ruolo di
decisioni di recourse: il loro valore pu\`o dipendere dall'informazione
rivelata dopo la decisione iniziale.

In termini della notazione generale della programmazione stocastica,

\[
x
=
\begin{pmatrix}
x_1\\
x_2\\
x_3\\
x_4
\end{pmatrix}
\]

\`e quindi la decisione di primo stadio, mentre le variabili

\[
g_1,\ldots,g_4,
\qquad
u_1,u_2,u_3
\]

costituiscono le componenti economicamente rilevanti della decisione di
recourse $y(\xi)$.

\section{Il caso Alcoa: mercato elettrico e struttura della decisione}

\label{sec:cap14-alcoa-mercato}

Il secondo caso guida considera il problema di approvvigionamento elettrico
di un grande consumatore industriale. Il riferimento metodologico \`e il
modello di grande consumatore sviluppato da Conejo, Carri\'on e
Morales.\footnote{A. J. Conejo, M. Carri\'on e J. M. Morales,
\emph{Decision Making Under Uncertainty in Electricity Markets},
Springer, 2010, Capitolo 9.}

La produzione di alluminio costituisce un contesto naturale per questo tipo
di problema, poich\'e il consumo di energia elettrica rappresenta una
componente rilevante del processo produttivo e del costo industriale.
Utilizzeremo pertanto Alcoa San Cipri\'an come riferimento economico del
caso.

Il modello che verr\`a costruito non intende tuttavia ricostruire i contratti
effettivamente sottoscritti da Alcoa, n\'e riprodurre integralmente il
modello di Conejo et al. L'obiettivo \`e isolare una struttura sufficientemente
realistica da rendere economicamente significativa la decisione, ma
sufficientemente semplice da poter essere studiata come problema lineare di
programmazione stocastica a due stadi.

\subsection{Il mercato dell'energia elettrica}

L'energia elettrica presenta caratteristiche che rendono il suo mercato
diverso da quello di molte altre commodity. Produzione e consumo devono
rimanere continuamente bilanciati e le possibilit\`a di immagazzinamento,
pur crescenti, restano limitate rispetto alla dimensione complessiva del
sistema. Il valore dell'energia dipende quindi non soltanto dalla quantit\`a,
ma anche dal periodo nel quale essa viene resa disponibile.

Nel mercato all'ingrosso operano diversi soggetti. I produttori immettono
energia nel sistema; grossisti e trader acquistano e vendono energia e
strumenti collegati; societ\`a di vendita e grandi consumatori acquistano
energia per soddisfare il proprio fabbisogno. Gli operatori di mercato
organizzano le piattaforme di negoziazione, mentre il gestore della rete di
trasmissione assicura l'equilibrio fisico del sistema.

Le transazioni possono essere concluse su mercati organizzati oppure
direttamente tra controparti. In quest'ultimo caso si parla di negoziazione
\emph{over the counter} (OTC). Un contratto OTC pu\`o essere costruito sulle
esigenze delle parti specificando quantit\`a, periodo di fornitura, profilo
temporale e regola di prezzo. I contratti bilaterali di approvvigionamento
costituiscono un esempio importante di questa forma di negoziazione.

Su un orizzonte pi\`u lungo operano anche i mercati dei derivati
sull'elettricit\`a. Tra gli strumenti pi\`u diffusi vi sono i futures, cio\`e
contratti standardizzati negoziati su mercati organizzati il cui valore
dipende dal prezzo futuro dell'energia. Essi permettono a produttori,
intermediari e consumatori di ridurre l'esposizione alle variazioni dei
prezzi che si realizzeranno successivamente sul mercato spot.

Il termine \emph{spot} comprende invece le negoziazioni prossime al momento
della fornitura fisica. Nel mercato \emph{day-ahead} vengono scambiate
quantit\`a di energia per il giorno successivo. I mercati
\emph{intraday} consentono di modificare successivamente le posizioni,
avvicinandole alle condizioni di produzione e consumo che diventano
progressivamente osservabili.

A ridosso della consegna interviene infine il meccanismo di bilanciamento,
attraverso il quale il gestore della rete corregge gli scostamenti tra
immissioni e prelievi e mantiene l'equilibrio fisico del sistema. Il mercato
di bilanciamento svolge quindi una funzione diversa da quella dei mercati
utilizzati principalmente per l'approvvigionamento o per la copertura del
rischio di prezzo.

Per il problema che svilupperemo sono particolarmente rilevanti tre
possibilit\`a:

\begin{itemize}

\item stipulare anticipatamente contratti bilaterali OTC per assicurarsi una
parte dell'energia futura a condizioni definite;

\item utilizzare contratti futures per coprire finanziariamente
l'esposizione alle variazioni dei prezzi;

\item acquistare o vendere energia sul mercato spot quando ci si avvicina al
periodo di consumo.

\end{itemize}

La dimensione temporale \`e particolarmente importante. La domanda di
elettricit\`a varia durante la giornata e, insieme alle condizioni
dell'offerta, determina profili di prezzo differenti nelle diverse ore.

Per descrivere questa variazione senza lavorare necessariamente a frequenza
oraria, \`e utile raggruppare le ore della giornata in \emph{fasce di
carico}. Una classificazione semplice distingue:

\begin{itemize}

\item \emph{base}, per le ore caratterizzate da domanda relativamente
bassa;

\item \emph{shoulder}, per le ore con condizioni intermedie;

\item \emph{peak}, per le ore nelle quali la domanda sul sistema \`e
relativamente elevata.

\end{itemize}

I confini esatti di queste fasce dipendono dal mercato e dal prodotto
considerato: base, shoulder e peak devono quindi essere intese come
categorie economiche utili a rappresentare differenti condizioni
intragiornaliere, non come una suddivisione universale delle ventiquattro
ore.

La stessa struttura temporale compare nei prodotti di approvvigionamento.
Un profilo \emph{base} prevede una fornitura costante lungo tutti i periodi
compresi nell'intervallo contrattuale; un profilo \emph{peak} concentra
invece la fornitura nei periodi di maggiore carico. Contratti con profili
differenti permettono quindi a un consumatore di costruire una copertura
pi\`u vicina alla forma temporale del proprio fabbisogno.

Per un grande consumatore industriale il problema non consiste pertanto
soltanto nell'acquistare una certa quantit\`a complessiva di elettricit\`a.
Occorre decidere quanto approvvigionamento fissare anticipatamente, quale
profilo temporale coprire e quanta esposizione lasciare alle condizioni che
si manifesteranno successivamente sul mercato spot.

\subsection{Un grande consumatore nel mercato elettrico}

Un grande consumatore industriale deve garantire, periodo per periodo, la
disponibilit\`a dell'energia necessaria al processo produttivo. Nel caso di
un impianto per la produzione di alluminio, il costo dell'elettricit\`a
rappresenta inoltre una componente rilevante del costo industriale.

Il problema di approvvigionamento combina quindi due esigenze: assicurare la
copertura del fabbisogno fisico e limitare l'esposizione a prezzi futuri
sfavorevoli.

Nel caso Alcoa consideriamo tre strumenti: contratti bilaterali stipulati
anticipatamente, transazioni successive sul pool e autoproduzione. La
decisione assume pertanto una struttura naturale:


\[
\boxed{
\begin{array}{c}
\text{scelta dei tre contratti bilaterali}
\\[3pt]
\downarrow
\\[3pt]
\text{rivelazione della traiettoria dei prezzi spot}
\\[3pt]
\downarrow
\\[3pt]
\text{transazioni spot e autoproduzione}.
\end{array}
}
\]


La prima scelta deve essere effettuata prima di conoscere i prezzi futuri;
le decisioni successive possono invece adattarsi alle condizioni che si
realizzano. Dopo la rivelazione dei prezzi l'impresa pu\`o acquistare o vendere energia
sul mercato spot e utilizzare la propria capacit\`a di autoproduzione.
Quest'ultima \`e rappresentata mediante un costo marginale costante e un
limite di capacit\`a.

La rivelazione progressiva dell'informazione su pi\`u date verr\`a considerata nel
capitolo successivo attraverso una formulazione multistadio.

\subsection{Orizzonte temporale, periodi e consumo}

L'orizzonte di pianificazione \`e di quattro settimane. Ogni giornata viene
suddivisa in tre periodi di otto ore, associati alle fasce base, shoulder e
peak. Si ottengono quindi $84$ periodi operativi.

Indichiamo l'insieme dei periodi con $\mathcal T=\{1,\ldots,84\}$ e con

\[
\mathcal T_{\mathrm{shoulder}}
=
\{2,5,8,\ldots,83\}, \qquad \mathcal T_{\mathrm{peak}}
=
\{3,6,9,\ldots,84\}
\]

l'insieme dei periodi shoulder e peak. Ciascun periodo ha durata

\[
\Delta_t=8
\]

ore, ove il pedice $t$ identifica l'intervallo temporale di otto ore, non un
istante.

Consideriamo tre contratti bilaterali, con profilo rispettivamente base,
shoulder e peak. La struttura temporale e la copertura offerta dai tre contratti bilaterali
sono illustrate nella Tabella~\ref{tab:alcoa-periodi}.

\begin{table}[!ht]
\centering
\caption{Fasce temporali e copertura dei consumi dei contratti bilaterali}
\label{tab:alcoa-periodi}
\small
\begin{tabular}{|l|ccccccccccc|}
\hline
Periodo
& 1 & 2 & 3 & 4 & 5 & 6 & 7 & 8 & 9 & $\cdots$ & 84 \\
\hline
Fascia
& B & S & P
& B & S & P
& B & S & P
& $\cdots$ & P \\
\hline
Contratto base
& $\times$ & $\times$ & $\times$
& $\times$ & $\times$ & $\times$
& $\times$ & $\times$ & $\times$
& $\cdots$ & $\times$ \\
Contratto shoulder
& & $\times$ &
& & $\times$ &
& & $\times$ &
& $\cdots$ & \\
Contratto peak
& & & $\times$
& & & $\times$
& & & $\times$
& $\cdots$ & $\times$ \\
\hline
\end{tabular}
\end{table}

Indichiamo con

\[
C_t
\]

il consumo elettrico dello stabilimento nel periodo $t$, espresso in MWh.
Nel modello il profilo dei consumi

\[
C_1,\ldots,C_{84}
\]

\`e noto. L'incertezza riguarda invece il prezzo futuro dell'energia sul
mercato spot.


\subsection{La decisione preventiva: i contratti bilaterali}

Indichiamo con

\[
x_{\mathrm B},
\qquad
x_{\mathrm{Sh}},
\qquad
x_{\mathrm P}
\]

le potenze, espresse in MW, contrattate rispettivamente mediante il contratto
base, shoulder e peak. Le quantit\`a contrattate sono grandezze continue e vengono tutte decise nel primo stadio.

La decisione di primo stadio \`e quindi

\[
x
=
\begin{pmatrix}
x_{\mathrm B}\\
x_{\mathrm{Sh}}\\
x_{\mathrm P}
\end{pmatrix}.
\]

Il contratto base opera in tutti i periodi dell'orizzonte; il contratto
shoulder soltanto per $t\in\mathcal T_{\mathrm{shoulder}}$; il contratto
peak soltanto per $t\in\mathcal T_{\mathrm{peak}}$.

Indichiamo con

\[
c_{\mathrm B}^{B},
\qquad
c_{\mathrm{Sh}}^{B},
\qquad
c_{\mathrm P}^{B}
\]

i rispettivi prezzi contrattuali, espressi in euro/MWh, e con

\[
\bar x_{\mathrm B},
\qquad
\bar x_{\mathrm{Sh}},
\qquad
\bar x_{\mathrm P}
\]

le massime potenze contrattabili. Le decisioni soddisfano quindi

\[
0\leq x_{\mathrm B}\leq\bar x_{\mathrm B},
\qquad
0\leq x_{\mathrm{Sh}}\leq\bar x_{\mathrm{Sh}},
\qquad
0\leq x_{\mathrm P}\leq\bar x_{\mathrm P}.
\]

L'energia fornita dai contratti nel periodo $t$ \`e

\[
\Delta_t
\left[
x_{\mathrm B}
+
\1_{\{t\in\mathcal T_{\mathrm{shoulder}}\}}x_{\mathrm{Sh}}
+
\1_{\{t\in\mathcal T_{\mathrm{peak}}\}}x_{\mathrm P}
\right].
\]

Il costo complessivo della copertura contrattuale \`e

\[
c_{\mathrm B}^{B}x_{\mathrm B}
\sum_{t\in\mathcal T}\Delta_t
+
c_{\mathrm{Sh}}^{B}x_{\mathrm{Sh}}
\sum_{t\in\mathcal T_{\mathrm{shoulder}}}\Delta_t
+
c_{\mathrm P}^{B}x_{\mathrm P}
\sum_{t\in\mathcal T_{\mathrm{peak}}}\Delta_t.
\]


\subsection{Il prezzo spot e le decisioni di recourse}

Indichiamo con

\[
c_t^{\mathrm{sp}}
\]

il prezzo dell'energia sul mercato spot nel periodo $t$, espresso in
euro/MWh. Al momento della scelta di $x$ i valori futuri di
$c_t^{\mathrm{sp}}$ non sono noti.

Nel formalismo generale del capitolo poniamo

\[
\xi_t=c_t^{\mathrm{sp}},
\qquad
t\in\mathcal T,
\]

cosicch\'e il processo

\[
\{\xi_t\}_{t\in\mathcal T}
=
\{c_t^{\mathrm{sp}}\}_{t\in\mathcal T}
\]

rappresenta l'incertezza del caso.

Dopo la rivelazione dei prezzi, indichiamo con

\[
y_t^{\mathrm{sp}}(\xi)\in\R
\]

la posizione netta sul mercato spot nel periodo $t$. Adottiamo la convenzione

\[
y_t^{\mathrm{sp}}(\xi)>0
\quad\text{per un acquisto},
\qquad
y_t^{\mathrm{sp}}(\xi)<0
\quad\text{per una vendita}.
\]

Indichiamo inoltre con

\[
y_t^G(\xi)\geq0
\]

l'energia autoprodotta nel periodo $t$. La capacit\`a disponibile impone

\[
0
\leq
y_t^G(\xi)
\leq
\bar y_t^G,
\]

mentre

\[
c^G
\]

indica il costo marginale unitario dell'autoproduzione.

Le quantit\`a

\[
y_t^{\mathrm{sp}}(\xi),
\qquad
y_t^G(\xi),
\qquad
t\in\mathcal T,
\]

costituiscono le decisioni operative di secondo stadio.

\subsection{Il bilancio energetico}

In ogni periodo $t\in\mathcal T$ l'energia disponibile deve essere sufficiente
a soddisfare il consumo $C_t$. Come gi\`a indicato, la componente fornita dai contratti bilaterali \`e

\[
\Delta_t
\left[
x_{\mathrm B}
+
\1_{\{t\in\mathcal T_{\mathrm{shoulder}}\}}x_{\mathrm{Sh}}
+
\1_{\{t\in\mathcal T_{\mathrm{peak}}\}}x_{\mathrm P}
\right].
\]

A questa si aggiungono l'autoproduzione $y_t^G(\xi)$ e la posizione netta sul
mercato spot $y_t^{\mathrm{sp}}(\xi)$. Il bilancio energetico \`e quindi

\[
\boxed{
\Delta_t
\left[
x_{\mathrm B}
+
\1_{\{t\in\mathcal T_{\mathrm{shoulder}}\}}x_{\mathrm{Sh}}
+
\1_{\{t\in\mathcal T_{\mathrm{peak}}\}}x_{\mathrm P}
\right]
+
y_t^G(\xi)
+
y_t^{\mathrm{sp}}(\xi)
=
C_t,
\qquad
t\in\mathcal T.
}
\]

Poich\'e $y_t^{\mathrm{sp}}(\xi)$ \`e libera in segno, lo stesso vincolo
permette sia un acquisto sia una vendita sul mercato spot. Se contratti e
autoproduzione non coprono il consumo, deve risultare
$y_t^{\mathrm{sp}}(\xi)>0$; se generano un'eccedenza, pu\`o risultare
$y_t^{\mathrm{sp}}(\xi)<0$.

Il bilancio collega quindi direttamente la decisione preventiva $x$ alle
decisioni di recourse: diversi contratti bilaterali modificano le
quantit\`a che dovranno essere acquistate, vendute o autoprodotte dopo
l'osservazione dei prezzi.


\subsection{Struttura dei costi e natura della decisione}

La componente di costo determinata nel primo stadio \`e quella relativa ai
contratti bilaterali. Poich\'e il contratto base opera in $84$ periodi,
mentre i contratti shoulder e peak operano ciascuno in $28$ periodi, e
$\Delta_t=8$, il costo contrattuale \`e

\[
C^B(x)
=
8
\left[
84\,c_{\mathrm B}^{B}x_{\mathrm B}
+
28\,c_{\mathrm{Sh}}^{B}x_{\mathrm{Sh}}
+
28\,c_{\mathrm P}^{B}x_{\mathrm P}
\right].
\]

Dopo la rivelazione dei prezzi spot, la componente gestionale (trading sul mercato spot e autoproduzione) del costo \`e

\[
C^R(y(\xi),\xi)
=
\sum_{t=1}^{84}
\left[
c_t^{\mathrm{sp}}y_t^{\mathrm{sp}}(\xi)
+
c^G y_t^G(\xi)
\right].
\]

Quando $y_t^{\mathrm{sp}}(\xi)<0$, il termine
$c_t^{\mathrm{sp}}y_t^{\mathrm{sp}}(\xi)$ \`e negativo e rappresenta il
ricavo derivante dalla vendita di energia sul mercato spot. La decisione di vendere energia  non \`e affatto irrealistica. Essa pu\`o verificarsi in diversi casi quali, ad esempio se in certo tempo $t$ si verificasse,
\[
c_t^{\mathrm{sp}}(\xi) > c^G.
\]
converrebbe autogenerare energia e vendere l'eventuale eccesso non consumato sul mercato spot.

Il costo complessivo associato alla decisione iniziale $x$ e alla
realizzazione dell'incertezza \`e pertanto

\[
\Psi(x,y(\xi),\xi)
=
C^B(x)
+
C^R(y(\xi),\xi).
\]

La struttura del caso pu\`o essere sintetizzata come

\[
\boxed{
x
\quad\longrightarrow\quad
\{c_t^{\mathrm{sp}}\}_{t\in\mathcal T}
\quad\longrightarrow\quad
y(\xi).
}
\]

La decisione $x$ determina anticipatamente il livello e il profilo temporale
della copertura contrattuale. Le decisioni di secondo stadio permettono
invece di adattare l'approvvigionamento ai prezzi effettivamente osservati.

Una maggiore copertura contrattuale riduce il fabbisogno da soddisfare
successivamente sul mercato spot, ma vincola anticipatamente l'impresa a
prezzi prefissati. Una copertura minore aumenta la flessibilit\`a, ma lascia
una quota maggiore del consumo esposta all'incertezza dei prezzi.

Nelle sezioni successive questa struttura verr\`a completata introducendo
gli scenari e le loro probabilit\`a. Il costo atteso costituir\`a il criterio
di riferimento; per il caso Alcoa verr\`a inoltre considerato il
Conditional Value-at-Risk del costo complessivo, cos\`i da distinguere la
soluzione risk-neutral da quella di un consumatore avverso al rischio.

\section{Stadi decisionali e struttura dell'informazione}

\label{sec:cap14-stadi-informazione}

La programmazione stocastica richiede di specificare non soltanto quali
quantit\`a siano incerte, ma anche quale informazione sia disponibile nel
momento in cui ciascuna decisione viene assunta.

Il Capitolo~5 ha rappresentato il flusso dell'informazione mediante una
filtrazione

\[
\{\F_t\}_{t\in\mathcal T},
\]

con

\[
\F_s\subseteq\F_t
\qquad
\text{per }s\leq t.
\]

Riprendiamo ora questa struttura per stabilire quali decisioni siano
ammissibili prima e dopo la rivelazione dell'incertezza.


\subsection{Periodi temporali e stadi decisionali}

Un \emph{periodo temporale} identifica un intervallo dell'orizzonte
economico. Uno \emph{stadio decisionale} identifica invece un momento nel
quale viene assunta una decisione sulla base dell'informazione allora
disponibile.

I due concetti non coincidono.

Il caso Alcoa, per esempio, contiene $84$ periodi operativi, ma nel modello
del presente capitolo soltanto due stadi decisionali. Nel primo vengono
scelti i contratti bilaterali; nel secondo vengono determinate le decisioni
di spot e autoproduzione dopo la rivelazione dell'incertezza rilevante.

In generale, pi\`u decisioni appartengono allo stesso stadio quando vengono
assunte sulla base della stessa informazione. Un nuovo stadio si apre quando
l'arrivo di nuova informazione consente di differenziare le decisioni in
funzione di quanto osservato.

Per rappresentare i due livelli informativi indichiamo con

\[
\F^{(1)}
\qquad\text{e}\qquad
\F^{(2)}
\]

le sigma-algebre disponibili rispettivamente al primo e al secondo stadio,
con

\[
\F^{(1)}
\subseteq
\F^{(2)}
\subseteq
\F.
\]

Esse possono essere interpretate come due livelli della filtrazione
corrispondenti ai momenti nei quali vengono assunte le decisioni.


\subsection{Informazione disponibile e decisioni ammissibili}

La decisione di primo stadio

\[
x
\]

deve utilizzare soltanto l'informazione disponibile prima della rivelazione
dell'incertezza. Formalmente, le sue componenti devono essere
$\F^{(1)}$-misurabili.

Quando l'informazione iniziale \`e fissata nel momento in cui il problema
viene risolto, $x$ viene quindi rappresentato semplicemente come un vettore
deterministico.

La decisione di secondo stadio

\[
y(\xi)
\]

pu\`o invece utilizzare l'informazione successivamente acquisita e deve
essere $\F^{(2)}$-misurabile.

La misurabilit\`a esprime un vincolo economico essenziale: una decisione pu\`o
dipendere da ci\`o che \`e gi\`a stato osservato, ma non da informazioni che
non sono ancora disponibili.

Non si deve quindi scrivere

\[
x\in\F^{(1)}
\qquad\text{o}\qquad
y(\xi)\in\F^{(2)}.
\]

Le sigma-algebre sono insiemi di eventi; una decisione \`e invece una
variabile o un vettore di variabili. La formulazione corretta richiede che
la decisione sia \emph{misurabile rispetto} alla sigma-algebra informativa
corrispondente.


\subsection{Rivelazione dell'incertezza nel modello a due stadi}

Il processo stocastico

\[
\{\xi_t\}_{t\in\mathcal T}
\]

rappresenta l'evoluzione dei dati esogeni incerti. L'informazione generata
dalla sua traiettoria sull'orizzonte pu\`o essere indicata con

\[
\sigma\left(\xi_t:t\in\mathcal T\right).
\]

Nel modello a due stadi assumiamo che l'informazione rilevante per il
recourse sia disponibile prima della decisione di secondo stadio. Possiamo
quindi rappresentare il secondo livello informativo come

\[
\F^{(2)}
=
\F^{(1)}
\vee
\sigma\left(\xi_t:t\in\mathcal T\right),
\]

dove $\vee$ indica la pi\`u piccola sigma-algebra contenente entrambe le
fonti di informazione.

La sequenza

\[
x
\quad\longrightarrow\quad
\xi
\quad\longrightarrow\quad
y(\xi)
\]

acquista cos\`i un significato preciso:

\[
\begin{array}{ccl}
x
&:&
\F^{(1)}\text{-misurabile},
\\[3pt]
\xi
&:&
\text{incertezza esogena rivelata tra i due stadi},
\\[3pt]
y(\xi)
&:&
\F^{(2)}\text{-misurabile}.
\end{array}
\]

Questa struttura implica che la decisione $x$ non possa essere modificata
in funzione della realizzazione futura di $\xi$, mentre il recourse pu\`o
differire tra realizzazioni diverse.

Quando nella sezione successiva l'incertezza verr\`a rappresentata mediante
scenari, questa distinzione avr\`a una conseguenza immediata: la decisione
$x$ sar\`a unica per tutti gli scenari, mentre a scenari differenti potranno
corrispondere decisioni di secondo stadio differenti.


\subsection{Valutazione condizionata all'informazione corrente}

La stessa struttura informativa entra nella valutazione economica delle
decisioni.

Sia

\[
\Psi(x,y(\xi),\xi)
\]

un criterio integrabile associato alla decisione iniziale, al recourse e
alla realizzazione dell'incertezza. Dal punto di vista del decisore al primo
stadio, la quantit\`a rilevante \`e

\[
\E
\left[
\Psi(x,y(\xi),\xi)
\mid
\F^{(1)}
\right].
\]

Il valore atteso condizionato utilizza quindi esattamente la nozione
introdotta nei Capitoli~3 e~5: valuta le conseguenze future utilizzando
soltanto l'informazione disponibile nel momento della decisione.

Quando l'informazione iniziale \`e fissata e non contiene ulteriore
incertezza rilevante, il criterio si riduce al consueto valore atteso

\[
\E
\left[
\Psi(x,y(\xi),\xi)
\right].
\]

La formulazione a scenari che verr\`a introdotta nel seguito fornir\`a una
rappresentazione finita di questo valore atteso.


\subsection{La struttura informativa nei due casi guida}

Nel caso Alcoa la decisione

\[
x
=
\begin{pmatrix}
x_{\mathrm B}\\
x_{\mathrm{Sh}}\\
x_{\mathrm P}
\end{pmatrix}
\]

viene assunta prima di conoscere la traiettoria dei prezzi spot. Poich\'e

\[
\xi_t=c_t^{\mathrm{sp}},
\]

nel modello a due stadi il secondo livello informativo contiene

\[
\sigma
\left(
c_t^{\mathrm{sp}}:
t\in\mathcal T
\right).
\]

Le decisioni $y_t^{\mathrm{sp}}(\xi)$ e $y_t^G(\xi)$ possono quindi dipendere
dalla traiettoria dei prezzi rivelata, mentre i tre contratti iniziali non
possono essere modificati.

Nel caso SVB--ALM l'incertezza esogena viene invece rappresentata mediante
il processo di stato

\[
\{Z_t\}.
\]

Quando $Z_t$ segue una catena di Markov rispetto alla filtrazione considerata,

\[
\Prob
\left(
Z_{t+1}=z_j
\mid
\F_t
\right)
=
\Prob
\left(
Z_{t+1}=z_j
\mid
Z_t
\right).
\]

Se lo stato corrente \`e noto ed \`e pari a $z_i$, si ha quindi

\[
\Prob
\left(
Z_{t+1}=z_j
\mid
Z_t=z_i
\right)
=
p_{ij}.
\]

L'informazione contenuta nell'intera storia passata viene cos\`i sintetizzata,
per la distribuzione dello stato successivo, dallo stato corrente. La riga
$i$ della matrice di transizione determina le probabilit\`a dei possibili
stati futuri.

I due casi utilizzano quindi modelli probabilistici differenti, ma rispettano
lo stesso principio informativo: la decisione iniziale utilizza soltanto
l'informazione corrente, mentre il recourse pu\`o reagire all'informazione
resa disponibile tra i due stadi.